\chapter{Introduction}


\section{Motivation}
The term structure of interest rates, commonly referred to as the yield curve, is a fundamental object in finance and macroeconomics. It summarizes market expectations about future interest rates, inflation, and economic conditions, and serves as a key input for asset pricing, portfolio allocation, and risk management. For policymakers, the yield curve provides valuable information about the stance of monetary policy and expectations about future economic activity, while for financial institutions it is central to valuation, hedging, and profitability analysis.

A large body of literature has proposed models to describe and forecast the yield curve. Parametric models such as the Nelson–Siegel and Nelson–Siegel–Svensson specifications offer a parsimonious and intuitive representation of the yield curve through a small number of latent factors. In parallel, stochastic interest rate models, such as the Hull–White model, are widely used in practice due to their analytical tractability and suitability for derivative pricing. Despite their popularity, these frameworks often rely primarily on yield curve information and treat macroeconomic conditions only implicitly, if at all.

Recent empirical evidence suggests that macroeconomic variables—such as inflation, output growth, and monetary policy rates—play a significant role in shaping the dynamics of the yield curve. Ignoring these factors may limit traditional models' ability to explain yield movements and produce accurate forecasts, particularly during periods of economic stress or structural change. This observation motivates integrating macroeconomic information into yield curve modeling frameworks.

\section{Research Objectives and Questions}
The primary objective of this thesis is to examine whether yield curve models that explicitly incorporate macroeconomic factors provide improved explanatory power and forecasting performance compared to standard stochastic interest rate models. To address this objective, the study focuses on the following research questions:

- To what extent do macroeconomic variables explain the dynamics of the latent factors of the yield curve?
- Does a macro-augmented econometric yield curve model outperform the Hull–White one-factor model in terms of in-sample fit and out-of-sample forecasting accuracy?
- How do the two modeling approaches differ in their economic interpretability and response to macroeconomic shocks?

By answering these questions, the thesis aims to assess the practical relevance of macro-financial integration in yield curve modeling.


\section{Contribution of the Thesis}
This thesis contributes to the existing literature in three main ways. First, it extends a widely used yield curve modeling framework by linking its latent factors to observable macroeconomic variables using econometric methods. Second, it provides a direct empirical comparison between a macro-augmented yield curve model and the Hull–White stochastic interest rate model, which remains a standard benchmark in financial practice. Third, the study offers insights into the economic interpretation of yield curve movements by explicitly relating them to macroeconomic conditions.

The results of this research are relevant for both academics and practitioners, particularly in the context of monetary policy analysis, interest rate forecasting, and financial risk management.

